\lab{7}{Булев индекс}
\task
Построить собственный бинарный обратный индекс и прямой индекс заголовков
и URL. Описать побайтовый формат, внутреннее представление, сортировку,
количество и среднюю длину термов, ограничения скорости и масштабирования.

\method
Корпус читается последовательно. Каждому документу назначается
\(docID\) по порядку в потоке. Токенизация и стемминг создают записи
«терм, документ, частота». Буфер ограничен параметром \code{memory_mib}.
После заполнения он сортируется по терму и \(docID\), одинаковые пары
объединяются суммированием частот, результат записывается во временный
run-файл. Отсортированные файлы попарно сливаются до единого потока.
Это вариант блочной индексации с внешней сортировкой.

Для массивов, сортировки и операций со списками используются собственные
структуры и функции. Файл записывается явно по полям, а не выгрузкой
структуры C++: выравнивание компилятора не определяет формат.
Все многобайтные целые имеют порядок little-endian.

\begin{longtable}{rlp{0.53\textwidth}}
\caption{Заголовок булевого индекса IRIDX001, 96 байт.}\\
\toprule Смещение & Размер & Поле\\\midrule\endfirsthead
\toprule Смещение & Размер & Поле\\\midrule\endhead
0 & 8 & Сигнатура \code{IRIDX001}\\
8 & 4 & Версия формата: 1\\
12 & 4 & Флаги нормализации\\
16 & 8 & Число документов\\
24 & 8 & Число термов\\
32 & 8 & Число токенов\\
40 & 8 & Объём входного текста\\
48 & 8 & Начало словаря\\
56 & 8 & Начало строк термов\\
64 & 8 & Начало постингов\\
72 & 8 & Начало смещений документов\\
80 & 8 & Начало метаданных\\
88 & 8 & Размер файла\\\bottomrule
\end{longtable}
Запись терма занимает 32 байта: смещение текста (8), смещение списка
постингов (8), суммарная частота (8), длина UTF-8 в байтах (4),
число документов (4). Смещения терма отсчитываются от соответствующих
секций. Словарь отсортирован лексикографически; поиск выполняется
бинарным поиском. Список постингов содержит возрастающие \code{uint32}
идентификаторы, по 4 байта на документ. Частота хранится агрегатом
в словаре, отдельной \(tf\) в булевом постинге нет.

Прямой индекс содержит 8-байтовые смещения записей документов.
В каждой записи идут две 4-байтовые длины, затем URL и заголовок.
Отделение прямого индекса позволяет извлекать метаданные только для
показываемых результатов. Флаги и версия создают основу для расширения,
но произвольный новый формат старый читатель не поддержит автоматически.

\journal
Построение разделено на накопление блока, сброс, слияние и сборку секций.
Внутри блока одинаковые записи схлопываются до записи на диск.
Для защиты данных добавлены проверки размеров, порядка постингов,
диапазона идентификаторов и границ секций. Корпус читается потоком,
поэтому его полный размер не определяет объём оперативной памяти.
При этом лимит блока не является строгим лимитом RSS всего процесса:
дополнительная память нужна строкам, буферам и служебным структурам.

\testing
Тест \code{ExternalSortAndMetadata} строит один корпус с буфером
128 байт и 1 MiB и сравнивает файлы побайтово. Проверяются суммарная
частота, число документов терма, отсутствие неизвестного слова и
заголовок Unicode. Второй тест проверяет NUL и обрезанные данные.
Оба теста пройдены. Эти тесты проверяют независимость
результата от разбиения на блоки, а не скорость на полном корпусе.

\results
\begin{table}[H]\centering
\begin{tabular}{lr}\toprule
Показатель & Значение\\\midrule
Документов & 1\,679\,971\\
Токенов & 2\,483\,697\,508\\
Термов после стемминга & 19\,787\,654\\
Средняя длина терма, кодовых точек & 13,2517\\
Размер булевого индекса, байт & 3\,618\,409\,348\\\bottomrule
\end{tabular}
\caption{Статистика готового булевого индекса, повторно прочитанная из файла.}
\end{table}
Средняя длина уникального терма существенно больше средней длины
вхождения токена: частые короткие служебные слова в первой величине
имеют вес один, а длинные редкие идентификаторы образуют большую часть
словаря. Стемминг укорачивает некоторые слова, но не устраняет этот эффект.

Время полного построения булевого индекса отдельно не измерялось.
Поэтому его скорость не выводится из времени сборки координатного индекса.
Параметр 64 MiB в конфигурации - размер рабочего бюджета блока,
а не измеренное пиковое потребление памяти. При росте корпуса в 10,
100 и 1000 раз растут число временных блоков и объём слияния.
Дисковая пропускная способность и число проходов становятся ограничением;
хранение \(docID\) в 32 битах дополнительно ограничивает число документов.

\conclusions
Создан переносимый бинарный формат с прямым и обратным индексом.
Внешняя сортировка позволяет строить его без загрузки всех документов
в память. Слабое место - повторное чтение временных файлов при
попарном слиянии. Многопутевое слияние и уменьшение числа проходов
могут снизить дисковую нагрузку. Для полноценной оценки скорости
необходимо сохранять журнал построения вместе со статистикой машины.
